\thispagestyle{plain}

\chapter{Data Loading and Inspection}
\label{cp:dli}

\vspace{2 \baselineskip}

The initial data inspection using df.info() reveals that the dataset consists of 48,842 entries, with 6 numerical columns and 9 categorical (object-type) columns.


\vspace{1 \baselineskip}
\begin{center}
\begin{tabular}{||c | c||} 
 \hline
 Qualitative variables & Quantitative variables \\ [1 ex] 
 \hline\hline
 Workclass & Age \\ [1ex] 
 \hline
 Marital-status & Fnlwgt \\ [1ex] 
 \hline
 Education & Educational-num \\ [1ex] 
 \hline
 Occupation & Capital-gain  \\ [1ex] 
 \hline
 Relationship & Capital-loss  \\ [1ex] 
 \hline
 Race & Hours-per-week  \\ [1ex] 
 \hline
 Gender & -  \\ [1ex] 
 \hline
 Native-Country & -  \\ [1ex] 
 \hline
 Income & -  \\ [1ex] 
 \hline
\end{tabular}
\end{center}


Notably, pandas reports zero non-null values across all columns. This is highly unusual and suggests that missing values might be encoded with a special character ('?', 'N/A') rather than the standard NaN format. We will investigate this hypothesis next.\\
The value\_counts output for the workclass column confirms our earlier suspicion. It reveals a ? category, which indicates that missing values are not encoded in a way that pandas can recognize automatically. To rectify this, we will replace all occurrences of the ? character with np.nan across the entire dataset. This will allow us to use standard pandas functions to analyze and handle these missing values. \\

The count of these missing values reveals the following:
\begin{itemize}
    \item workclass: 2799 missing values

    \item occupation: 2809 missing values

    \item native-country: 857 missing values

\end{itemize}

This is a significant amount of missing data (2799 out of 48 842 rows is about 5.7\% for workclass). Simply deleting these rows would cause a significant loss of information.\\

Our Exploratory Data Analysis (EDA) phase, which begins now, will need to determine the best strategy for handling these NaN values.